\SourceLeaf{534}gefunden und hieraus
\[
\begin{aligned}
[2]&=12(18)+8(57)-5(239)\\
{}[5]&=7(18)+5(57)-3(239)\\
{}[13]&=9(18)+6(57)-4(239)
\end{aligned}
\]
ferner mit Hülfe der Tafeln
\[
\begin{aligned}
(268)&=-2[5]+2[13]-[17]\quad\raisebox{.4ex}{$\pdot$}\\
(38)&=-[5]\phantom{+2[13]}+2[17]
\end{aligned}
\]
und hieraus durch Elimination von $[17]$ und Einsetzen der zuvor erhaltenen Werthe von $[5]$, $[13]$
\[
(38)+2(268)=(18)-(57)-(239)
\]
Die Elimination von $(18)$ hat dann die neue Bestimmung ergeben
\[
\begin{aligned}
[2]&=12(38)+20(57)+7(239)+24(268)\\
{}[5]&=7(38)+12(57)+4(239)+14(268)\\
{}[13]&=9(38)+15(57)+5(239)+18(268)\\
{}[17]&=4(38)+6(57)+2(239)+7(268)
\end{aligned}
\]
\par Nach folgeweiser Anwendung der Cotangenten 117, 327, 882, 18543, 307, 278, 378, 829, 993, 2943, 447, 606, 931, 1143, 1772, 6118, 34208, 44179, 85353, 485298, 17772, 9466, 330182, 5257, 114669, 12943 sind endlich $[2][5]\mathinner{.\,.\,.}[61]$ durch $(5257)$, $(9466)\mathinner{.\,.\,.}(485298)$ ausgedrückt und deren Coëfficienten in den folgenden Spalten zusammengestellt:
\begingroup\fontsize{9.5}{12}\selectfont\setlength{\tabcolsep}{2.3pt}
\begin{center}
\begin{tabular}{r|r|r|r|r|r|r|r|r|r}
 & $5257$ & $9466$ & $12943$ & $34208$ & $44179$ & $85353$ & $114669$ & $330182$ & $485298$ \\\hline
$2$ & $+2805$ & $-398$ & $+1950$ & $+1850$ & $+2021$ & $+2097$ & $+1484$ & $+1389$ & $+808$ \\
$5$ & $+1656$ & $-235$ & $+1151$ & $+1092$ & $+1193$ & $+1238$ & $+876$ & $+820$ & $+477$ \\
$13$ & $+2100$ & $-298$ & $+1460$ & $+1385$ & $+1513$ & $+1570$ & $+1111$ & $+1040$ & $+605$ \\
$17$ & $+875$ & $-124$ & $+608$ & $+577$ & $+630$ & $+654$ & $+463$ & $+433$ & $+252$ \\
$29$ & $+1359$ & $-193$ & $+945$ & $+896$ & $+979$ & $+1016$ & $+719$ & $+673$ & $+391$ \\
$37$ & $+590$ & $-84$ & $+410$ & $+389$ & $+425$ & $+441$ & $+312$ & $+292$ & $+170$ \\
$41$ & $+2410$ & $-342$ & $+1675$ & $+1589$ & $+1736$ & $+1802$ & $+1275$ & $+1193$ & $+694$ \\
$53$ & $+994$ & $-141$ & $+691$ & $+655$ & $+716$ & $+743$ & $+526$ & $+492$ & $+286$ \\
$61$ & $+2481$ & $-352$ & $+1725$ & $+1637$ & $+1788$ & $+1855$ & $+1313$ & $+1229$ & $+715$ \\
\end{tabular}
\end{center}\endgroup
\par Von der Richtigkeit dieser Gleichungen, welche zur Bestimmung von $[2][5]\dotsfour[61]$ dienen können, überzeugt man sich unmittelbar durch die aus obigen Tafeln sich ergebenden Zerlegungen
\begingroup\fontsize{10}{13}\selectfont\setlength{\arraycolsep}{1.3pt}
\[
\begin{array}{r@{\;}rrrrrrrrr}
(5257)={} & [2] & +2[5] & -[13] & +[17] & \pdot & \pdot & -[41] & \pdot & -[61] \\
(9466)={} & 2[2] & \pdot & \pdot & \pdot & -[29] & -3[37] & \pdot & \pdot & -[61] \\
(12943)={} & [2] & -4[5] & +3[13] & \pdot & \pdot & \pdot & \pdot & \pdot & -[61] \\
(34208)={} & 2[2] & -[5] & -2[13] & +[17] & +[29] & \pdot & \pdot & -2[53] & \pdot \\
(44179)={} & 3[2] & \pdot & -3[13] & -2[17] & -[29] & \pdot & \pdot & +[53] & \pdot \\
(85353)={} & -[2] & -[5] & +[13] & -[17] & \pdot & -[37] & +2[41] & -[53] & \pdot \\
(114669)={} & -3[2] & \pdot & \pdot & +[17] & \pdot & +[37] & \pdot & +2[53] & +2[61] \\
(330182)={} & -4[2] & +5[5] & +[13] & \pdot & +[29] & -[37] & -[41] & \pdot & +[61] \\
(485298)={} & -2[2] & -[5] & +4[13] & \pdot & -2[29] & +[37] & \pdot & +[53] & \pdot \\
\end{array}
\]
\endgroup
